\section{Radiation, neutrinos, and the supernova profile}
\label{sec:rad}

The second application profile, core-collapse supernovae, involves two
closures: one for neutrino transport, which is a moment closure with a long
history in radiation transport, and one for turbulence in a
spherically symmetric model, which is a Reynolds closure. Both are
recorded here because the first supplies the cleanest worked example of a
maximum-entropy closure and the second defines the benchmark target.

\subsection{The moment hierarchy of transport}

The specific intensity $I(\bx,\boldsymbol\Omega,t)$ obeys a linear kinetic
equation in which particles all move at speed $c$,
\[
  \frac1c\,\partial_tI+\boldsymbol\Omega\cdot\nabla I
  =-\sigma_tI+\frac{\sigma_s}{4\pi}\int I\,\dd\Omega'+S .
\]
Angular moments give the energy density $E=c^{-1}\int I\,\dd\Omega$, the
flux $\bF=\int\boldsymbol\Omega\,I\,\dd\Omega$ and the pressure tensor
$\mathbf P=c^{-1}\int\boldsymbol\Omega\boldsymbol\Omega\,I\,\dd\Omega$,
with
\begin{equation}
  \partial_tE+\nabla\cdot\bF=\text{sources},\qquad
  \frac1{c^2}\,\partial_t\bF+\nabla\cdot\mathbf P=-\frac{\sigma_t}{c}\,\bF .
\end{equation}
The pressure is unclosed. With $f=|\bF|/(cE)$ and
$\mathbf n=\bF/|\bF|$ the standard closures are written through an
\emph{Eddington factor} $\chi$,
\begin{equation}
  \mathbf P=\Bigl[\frac{1-\chi}{2}\,\mathbf I
     +\frac{3\chi-1}{2}\,\mathbf n\mathbf n\Bigr]E .
\end{equation}
Realizability requires $f\le1$ and $f^2\le\chi\le1$.

\subsection{The closures}

\begin{itemize}
  \item \textbf{Diffusion.} $\chi=\tfrac13$ and the flux equation in
    quasi-steady state: $\bF=-\tfrac{c}{3\sigma_t}\nabla E$. This is the
    Chapman--Enskog closure, valid when the mean free path $1/\sigma_t$ is
    short. In the opposite limit it gives $|\bF|>cE$: energy moving faster
    than light.
  \item \textbf{Flux-limited diffusion.} Keep the diffusion form with a
    limited coefficient,
    $\bF=-\tfrac{c\,\lambda(R)}{\sigma_t}\nabla E$ with
    $R=|\nabla E|/(\sigma_tE)$. The Levermore--Pomraning limiter
    \cite{LevermorePomraning81} is
    $\lambda(R)=(\coth R-1/R)/R$, which tends to $\tfrac13$ as $R\to0$ and
    to $1/R$ as $R\to\infty$, so that $|\bF|\to cE$. It is the transport
    analogue of the electron flux limiter of Section~\ref{sec:plasma}, with
    a derivation behind the limiter.
  \item \textbf{M1.} Evolve both $E$ and $\bF$ and close with a function
    $\chi(f)$.
  \item \textbf{Variable Eddington tensor.} Compute
    $\mathbf P/E$ from a separate, cheaper solution of the transport
    equation and use it to close the moment system.
\end{itemize}

\begin{example}[Two maximum-entropy Eddington factors]
\label{ex:m1}
Assume the intensity is isotropic in a frame moving with some velocity
$\beta c$, so that $I\propto(1-\beta\mu)^{-4}$ with $\mu$ the direction
cosine. This is how Levermore derived the closure \cite{Levermore84}, and
it coincides with maximizing the entropy of a photon gas subject to $E$
and $\bF$. Its moments give the closed form
\begin{equation}
  \chi(f)=\frac{3+4f^2}{5+2\sqrt{4-3f^2}} ,
  \label{eq:m1}
\end{equation}
with $\chi(0)=\tfrac13$ and $\chi(1)=1$. Maximizing the classical
(Boltzmann) entropy instead gives $I\propto e^{b\mu}$, with
$f=\coth b-1/b$ and $\chi=1-2f/b$ \cite{Minerbo78}. Both were checked by
quadrature: for example $\beta=0.4$ gives $f=0.5063$ and $\chi=0.4684$
from the integrals and from \eqref{eq:m1}.
\end{example}

The M1 closure is hyperbolic, realizable and exact in both limits, which
is everything Section~\ref{sec:moment} asked of a maximum-entropy closure.
Its known failure shows what the ansatz discards: for two beams crossing
head-on $f=0$ and the true $\chi$ is $1$, while \eqref{eq:m1} gives
$\tfrac13$. A function of $f$ alone cannot distinguish an isotropic field
from two opposed beams. Neutrino transport in supernova simulations has the
same structure with energy dependence and several species, and M1 is in
common use there.

\subsection{Turbulence in a one-dimensional supernova model}

A spherically symmetric model evolves angle averages. Let
$\avg{\cdot}_\Omega$ be the average over angles at fixed radius. Averaging
the three-dimensional equations and subtracting the one-dimensional
equations applied to the averaged state leaves
\begin{equation}
  C^\star_{\mathrm{SN}}[X]
  =\avg{F_{3\mathrm D}(X)}_\Omega-F_{1\mathrm D}\bigl(\avg{X}_\Omega\bigr),
\end{equation}
which is \eqref{eq:defect} for this projection. It decomposes into a
Reynolds stress (a turbulent pressure that helps support the shock), a
turbulent energy or enthalpy flux, production, dissipation and mixing. The
classical closure is mixing-length theory, the Prandtl model of
Section~\ref{sec:turb} with buoyancy as the driver, used in
one-dimensional explosion models \cite{CouchWarrenOConnor20}. What makes
this harder than the terrestrial case is that the flow is compressible,
passes through a moving shock, and never reaches a statistically steady
state, so history enters.
